\documentclass[12pt]{article}

\usepackage{amsmath,amssymb,amsthm}
\usepackage[margin=1in]{geometry}
\usepackage{xcolor}
%\usepackage{url}
%\usepackage{float}
%\usepackage{array}
\usepackage{comment}
\usepackage{booktabs}
\usepackage{graphicx}
\usepackage{tcolorbox}
%\usepackage{tikz}
%\usetikzlibrary{decorations.pathreplacing,arrows.meta,calc,patterns}
\usepackage{enumitem}
%\usepackage{csvsimple}
\usepackage{hyperref}
\usepackage[capitalize,sort]{cleveref}
%\usepackage{algorithm,algpseudocode}

% colors
\colorlet{textColor}{black}
\colorlet{bgColor}{white}
\colorlet{textRed}{red!50!black}
\colorlet{textGreen}{green!50!black}
\colorlet{textBlue}{blue!50!black}
\colorlet{textHeavy}{black}
\colorlet{dimColor}{textColor!50!bgColor}

\hypersetup{
colorlinks,linkcolor=textRed,citecolor=textRed,urlcolor=textBlue,
hypertexnames=false,bookmarksnumbered=true,final
}

% macros
\newcommand*{\Th}{^{\textrm{th}}}
\newcommand*{\abs}[1]{\lvert #1 \rvert}
\newcommand*{\defeq}{:=}
\DeclareMathOperator*{\E}{\mathbb{E}}
\DeclareMathOperator*{\Var}{Var}
\DeclareMathOperator*{\argmin}{argmin}
\DeclareMathOperator*{\argmax}{argmax}
\renewcommand*{\Pr}{\operatorname{\mathrm{P}}}

% theorems, lemmas, definitions
\theoremstyle{plain}
\newtheorem{theorem}{Theorem}
\newtheorem{lemma}[theorem]{Lemma}
\theoremstyle{definition}
\newtheorem{definition}{Definition}
\newtheorem{example}{Example}
\newtheorem{remark}{Remark}

% wrap URLs
\Urlmuskip=0mu plus 0.1mu
\makeatletter
\g@addto@macro{\UrlBreaks}{%
\do\/%
\do\a\do\b\do\c\do\d\do\e\do\f\do\g\do\h\do\i\do\j\do\k\do\l\do\m%
\do\n\do\o\do\p\do\q\do\r\do\s\do\t\do\u\do\v\do\w\do\x\do\y\do\z%
\do\A\do\B\do\C\do\D\do\E\do\F\do\G\do\H\do\I\do\J\do\K\do\L\do\M%
\do\N\do\O\do\P\do\Q\do\R\do\S\do\T\do\U\do\V\do\W\do\X\do\Y\do\Z%
\do\0\do\1\do\2\do\3\do\4\do\5\do\6\do\7\do\8\do\9%
}
\makeatother

% tighter lists
\newenvironment*{tightemize}{\begin{itemize}[noitemsep]}{\end{itemize}}%
\newenvironment*{tightenum}{\begin{enumerate}[noitemsep]}{\end{enumerate}}%

% customizations
\allowdisplaybreaks
%\defaultaddspace=0.0em  % spacing in tables

\newcommand*{\bibPrep}{
\phantomsection
\addcontentsline{toc}{section}{References}
\bibliographystyle{plainurl}
}


\graphicspath{{../../build/case-studies/3-reinsurance-plots/}}
\newcommand*{\Xbar}{\overline{X}}
\newcommand*{\Dcal}{\mathcal{D}}
\newcommand*{\Ncal}{\mathcal{N}}
\newcommand*{\Qhat}{\widehat{Q}}
\DeclareMathOperator{\Exp}{Exp}

\tcbuselibrary{theorems, skins, breakable}
\newtheorem{exercise}{Exercise}

\colorlet{mathExeBgColor}{blue!6!bgColor}
\colorlet{progExeBgColor}{orange!10!bgColor}
\colorlet{mathExeSec}{textColor!70!blue!40!bgColor}
\colorlet{progExeSec}{textColor!70!orange!45!bgColor}
\tcbset{
mathExe/.style={colback=mathExeBgColor, frame empty, colupper=textColor},
progExe/.style={colback=progExeBgColor, frame empty, colupper=textColor},
myres/.style={colback=textColor!3!bgColor, frame empty, colupper=textColor},
}
% \newtcbtheorem[number within=section]{exercise}{Exercise}{mytheoremstyle}{exe}

\title{IE 300 Case Study 3:\\Heavy-Tailed Distribution and Reinsurance Rate-Making}
\author{\empty}
\date{\empty}

\begin{document}

\maketitle
%\setlength{\parindent}{0pt}
\setlength{\parskip}{0.2em}

The purpose of this case study is to give a brief introduction to a heavy-tailed distribution
and its distinct behaviors in contrast with familiar light-tailed distributions in standard texts.
You will learn about QQ-plot, which is a popular tool for checking goodness-of-fit for
a particular statistical model. You will also work on a real-life application of heavy-tailed
distributions in reinsurance rate-making.

Reinsurance is a very important component of the global financial market.
It allows insurers to take on risks that they would otherwise not be able to.
Did you know that NASA buys insurance contracts for every rocket it launches
and every satellite and probe it sends to the outer space?
These equipments are so expensive that typical insurers would not be able to cover on their own.
Therefore, they can go to the reinsurance market, slide up the coverage and
transfer partial coverage to reinsurers that exceed their financial capabilities.
By the end of this case study, you will be able to learn basic principles of pricing a reinsurance contract.

\textbf{Learning objectives}:

\begin{tightenum}
\item Visualize the concepts in the Central Limit Theorem.
\item Identify cases where the Central Limit Theorem does not apply.
\item Reinforce the concept of cumulative distribution function.
\item Understand why and how QQ-plot works for the assessment of goodness-of-fit.
\item Reinforce the concepts of conditional probability and conditional expectation.
\item Apply basic integration technique to compute mean excess function.
\item Learn about behaviors of a heavy-tailed distribution.
\item Learn how to use order statistics to estimate quantiles and mean excess function.
\item Develop intuition behind point estimators.
\end{tightenum}

\newpage

\section{Central Limit Theorem}

This section is to provide visualization of central limit theorem which you should already be familiar with.
We provide examples for both discrete random variables and continuous random variables.

\subsection{Bernoulli Random Variables}
\label{sec:bernoulli}

Suppose we are given a potentially biased coin, i.e.,
a coin that may not have an equal chance of landing on a head or a tail.
We would like to estimate the probability of getting a head.

We can do so by counting the number of heads in a sequence of coin tosses.
Suppose we toss the coin $n$ times. For every $i$ from 1 to $n$, let $X_i$ be a random variable:
1 if the $i\Th$ coin shows heads and 0 otherwise.
Let $p$ the probability of a head and $q=1-p$ be the probability of a tail.
Then $X_i \sim \mathrm{Bernoulli}(p)$, i.e., $X_i$ is a Bernoulli random variable with parameter $p$.
Its probability mass function is given by
\[ \Pr(X_i = x) = \begin{cases}p & \text{ if } x = 1 \\ q & \text{ if } x = 0\end{cases}. \]

The total number of heads in $n$ tosses is given by
\[ S_n \defeq \sum_{i=1}^n X_i. \]
Then $S_n \sim \mathrm{Binomial}(n, p)$, i.e., $S_n$ is a binomial random variable
with parameters $n$ and $p$. Its probability mass function is given by
\[ \Pr(S_n = x) = \binom{n}{x}p^xq^{n-x}, \qquad x \in \{0, 1, \ldots, n\}. \]

The random variable
\[ \Xbar_n \defeq \frac{1}{n}S_n \]
is called the \emph{sample average}.
What are the mean and variance of $\Xbar_n$?
To find them, we will use the following result.

\begin{tcolorbox}[myres]
\begin{theorem}
\label[theorem]{thm:randvar-scaling}
For any real numbers $a$ and $b$, and random variable $Z$, we have
\begin{align*}
\E[aZ+b] &= a\E[Z]
& \Var[aZ+b] &= a^2\Var[Z]
\end{align*}
\end{theorem}
\end{tcolorbox}

\begin{figure}[!htb]
\centering
\includegraphics[scale=0.8]{binom1.pdf}
\caption{Probability mass function of $\Xbar_n$ for $p = 0.3$.}
\label{fig:binom1}
\end{figure}

Using \cref{thm:randvar-scaling}, we get
\[ \E[\Xbar_n] = \E[S_n]/n = p, \quad \text{and} \]
\[ \Var[\Xbar_n] = \Var\left[\frac{S_n}{n}\right] = \frac{\Var[S_n]}{n^2} = \frac{pq}{n}. \]
Since $\E[\Xbar_n] = p$, $\Xbar_n$ is an unbiased estimator of $p$.
Moreover, when $n$ is large, the variance of $\Xbar_n$ is small,
so $\Xbar_n$ is a good estimator of $p$.

Suppose that we have an unfair coin with $p = 0.3$.
\Cref{fig:binom1} shows the probability mass functions of the sample average $\Xbar_n$,
where the number of coin tosses $n \in \{1, 2, 3, 10, 20, 50\}$.
Since $p < 0.5$, we are more likely to see a smaller number of heads
than that of tails in any given $n$ tosses.
In general, the probability mass function of $\Xbar_n$ tends to skew towards to the right.
However, as one can see in the plots in \cref{fig:binom1},
the probability mass function becomes more and more symmetric as $n$ gets bigger and bigger.
This phenomenon arises for any $p \in (0, 1)$, no matter how extreme $p$ is.
Why is this happening? The answer is the Central Limit Theorem, which we have already learned in class.

The Central limit theorem tells us that $\Xbar_n$ is approximately normally distributed for large $n$,
i.e., $\Xbar_n$'s distribution is approximately $\Ncal(p, \frac{pq}{n})$.

Let us numerically verify it.
In \cref{fig:binom2}, for $p = 0.3$ and different values of $n$,
we plot the probability density function of $\Ncal(p, \frac{pq}{n})$ in dashed red,
and we plot $n$ times the probability mass function of $\Xbar_n$ in blue.
We can see that as $n$ increases, the two curves match up quite well.

\begin{figure}[!htb]
\centering
\includegraphics[scale=0.8]{binom2.pdf}
\caption{Probability mass function of $\Xbar_n$ for $p = 0.3$ vs normal distribution}
\label{fig:binom2}
\end{figure}

\subsection{Exponential Random Variables}

Let us now trying applying the Central limit theorem to the exponential distribution.
Let $X_1$, \ldots, $X_n$ be $n$ independent random variables such that for
all $i$ from $1$ to $n$, we have $X_i \sim \Exp(\lambda)$,
i.e., $X_i$ is exponentially distributed with parameter $\lambda$.
Recall that the probability density function of
an exponential random variable with parameter $\lambda$ is
\[ f(x) = \lambda e^{-\lambda x}, \quad x \ge 0, \]
and the cumulative distribution function is
\[ F(x) = 1 - e^{-\lambda x}, \quad x \ge 0, \]
the mean is $1/\lambda$, and the variance is $1/\lambda^2$.

Define $S_n$ and $\Xbar_n$ just like before:
\begin{align*}
S_n &\defeq \sum_{i=1}^n X_i,
& \Xbar_n &\defeq \frac{1}{n}S_n.
\end{align*}

Let us now find the mean and variance of $\Xbar_n$.
In \cref{sec:bernoulli}, we knew the distribution of $S_n$.
Here it's not so easy to see how $S_n$ might be distributed.
The following result can help us find $\Xbar_n$
without first finding out the distribution of $S_n$.

\begin{tcolorbox}[myres]
\begin{theorem}
\label[theorem]{thm:randvar-lincomb}
Let $Y_1, \ldots, Y_n$ be $n$ independent random variables.
Let $Z \defeq Y_1 + \ldots + Y_n$. Then
\begin{align*}
\E[Z] &= \sum_{i=1}^n \E[Y_i],
& \Var[Z] = \sum_{i=1}^n \Var[Y_i].
\end{align*}
\end{theorem}
\end{tcolorbox}

% \begin{exercise}{}{exe:exp:mean-var}
% \end{exercise}
\begin{tcolorbox}[mathExe]
\begin{exercise}
\label{exe:exp:mean-var}
Show that
\begin{enumerate}[label=(\alph*)]
\item $\E[\Xbar_n] = 1/\lambda$.
\item $\Var[\Xbar_n] = 1/(n\lambda^2)$.
\end{enumerate}
\end{exercise}
\end{tcolorbox}

The central limit theorem tells us that $\Xbar_n$ is approximately
normally distributed when $n$ is large.
By the results of \cref{exe:exp:mean-var}, we get that
$\Xbar_n$ is distributed roughly as $\Ncal(1/\lambda, 1/(n\lambda^2))$.

Let us try to verify the central limit theorem for $\Xbar_n$ numerically.

\begin{tcolorbox}[mathExe]
\begin{exercise}
\label{exe:exp:y-distr}
It can be shown that the probability density function of $S_n$ is given by
\[ f_{S_n}(x) = \frac{\lambda^n}{(n-1)!}x^{n-1}e^{-\lambda x}, \quad x \ge 0. \]
Define
\[ Y_n \defeq \frac{\Xbar_n - \E[\Xbar_n]}{\sqrt{\Var[\Xbar_n]}} = \frac{\lambda S_n - n}{\sqrt{n}}. \]
Show that the probability density function of $Y_n$ is
\[ f_{Y_n}(y) = \frac{\sqrt{n}}{(n-1)!}\left(n + y\sqrt{n}\right)^{n-1} \exp\left(-(n + y\sqrt{n})\right),
    \quad y \ge -\sqrt{n}. \]
\textcolor{mathExeSec}{Hint:
$F_{Y_n}(y) = \Pr(Y_n \le y) = \Pr(S_n \le (n + y\sqrt{n})/\lambda) = F_{S_n}((n + y\sqrt{n})/\lambda)$.}
\end{exercise}
\end{tcolorbox}

By the central limit theorem, $\Xbar_n$ should be distributed approximately $\Ncal(\E[\Xbar_n], \Var[\Xbar_n])$.
Thus, $Y_n$ should be distributed approximately $\Ncal(0, 1)$.
Since we found the true distribution of $Y_n$ in \cref{exe:exp:y-distr},
let us plot it against the standard normal distribution to see
how well the central limit theorem holds.

\begin{tcolorbox}[progExe]
\begin{exercise}
\label{exe:exp:y-plot}
Use \texttt{matplotlib.pyplot} to plot all of following functions in the same figure:
\begin{enumerate}
\item For $n \in \{1, 3, 10, 30, 100\}$, plot $f_{Y_n}(y)$ for $-4 \le y \le 4$.
    \\ \textcolor{progExeSec}{($Y_n$ is defined in \cref{exe:exp:y-distr}.
    $f_{Y_n}(y)$ is $0$ for $y < -\sqrt{n}$.)}
\item Plot the standard normal distribution's density function on the interval $[-4, 4]$ as a dashed line.
    \textcolor{progExeSec}{(You may use
    \texttt{\href{https://docs.scipy.org/doc/scipy/reference/generated/scipy.stats.norm.html}{%
scipy.stats.norm}.pdf(x)} to get the standard normal distribution's PDF at $x$.)}
\end{enumerate}
Ensure your curves are properly labeled with a legend.
\end{exercise}
\end{tcolorbox}

In general, we can conclude that Central Limit Theorem tells us that the distribution of
an average tends to be Normal, even when the distribution from which the average is computed is non-Normal.
However, there are certain cases in which the average deviates from Normal behavior.
When does the Central Limit Theorem not hold?

\subsection{Heavy-Tailed Distributions}

Heavy-tailed distributions are probability distributions with a heavier tail than the exponential.
We will see how the extremes produced by heavy-tailed distributions will corrupt the average
so that an asymptotic behavior different from the Normal behavior is obtained.
Formally, a random variable $X$ is said to have a heavy-tail if
\[ \lim_{x \to \infty} \frac{e^{-\lambda x}}{1 - F(x)} = 0, \qquad \text{for all } \lambda > 0. \]

Let us study the Pareto distribution. It has a single positive parameter $\alpha$,
called the \emph{Pareto index}, and its probability density function is given by
\[ f(x) = \alpha x^{-(\alpha + 1)}, \qquad x \ge 1. \]
The Pareto distribution is very important in reinsurance so we will study it closely.

\begin{tcolorbox}[mathExe]
\begin{exercise}
Show that
\begin{enumerate}[label=(\alph*)]
\item A Pareto random variable is heavy-tailed. \;
    \textcolor{mathExeSec}{(Hint: L'H\^opital's rule.)}
\item An exponential random variable with parameter $\alpha$ is not heavy-tailed.
\end{enumerate}
\end{exercise}

\begin{exercise}
Let $X \sim \mathrm{Pareto}(\alpha)$.
\begin{enumerate}[label=(\alph*)]
\item For what values of $\alpha$ does $\E[X]$ exist?
    When it exists, what is its value?
\item For what values of $\alpha$ do $\E[X]$ and $\Var[X]$ exist?
    When they exist, what is $\Var[X]$'s value.
\end{enumerate}
\end{exercise}

\begin{exercise}
The Pareto distribution is closely related to the exponential distribution.
Let $X \sim \mathrm{Pareto}(\alpha)$ and $Y = \ln(X)$.
Show that $Y \sim \Exp(1/\alpha)$.
\\ \textcolor{mathExeSec}{Hint: Use the same trick as in the hint of \cref{exe:exp:y-distr}.}
\end{exercise}
\end{tcolorbox}

Since the Pareto distribution may not have a finite mean or variance,
the Central limit theorem is not applicable.

\section{QQ plot}

The quantile-quantile plot, also known as a QQ-plot, is a tool to visually compare two distributions.
It can also be used to compare a dataset of samples to a known distribution.

\subsection{QQ plots of two distributions}

\begin{tcolorbox}[myres]
\begin{definition}[Quantile function]
\label{defn:quantile}
The \emph{quantile function} of a distribution is, intuitively,
the inverse of the cumulative distribution function.
Formally, for any $p \in [0, 1]$ and any random variable $X$, its quantile at $p$,
denoted as $Q_X(p)$, is the largest value $x$ such that $F_X(x) \le p$.
In other words, $Q_X(p) \ge x$ if and only if $F_X(x) \le p$.
\end{definition}
\end{tcolorbox}

\begin{tcolorbox}[myres]
\begin{example}[Exponential distribution]
Let $X \sim \Exp(\lambda)$. Then for any $p \in [0, 1]$,
\[ F_X(x) \le p
\iff 1 - e^{-\lambda x} \le p
\iff e^{\lambda x} \le \frac{1}{1-p}
\iff x \le \frac{1}{\lambda}\ln\left(\frac{1}{1-p}\right). \]
Thus,
\[ Q_X(p) = \frac{1}{\lambda}\ln\left(\frac{1}{1-p}\right). \]
\end{example}
\end{tcolorbox}

Let $X$ and $Y$ be two random variables.
The QQ plot of $Y$ w.r.t. $X$ is obtained by plotting the points
$(Q_X(p), Q_Y(p))$ for all $p \in [0, 1]$.

QQ plots can help us compare two distributions.
The idea is that if their distributions are similar,
then their quantile functions will be similar too,
and so, the plot will be very close to the line $y = x$.

QQ plots can also help us know if one random variable is a linear function of the other;
in that case, the QQ plot will be a straight line because of the following result.

\begin{tcolorbox}[myres]
\begin{theorem}
\label{thm:q-scale}
Let $X$ be a random variable and $Y \defeq aX + b$.
Then $Q_Y(p) = aQ_X(p) + b$.
\end{theorem}
\end{tcolorbox}
\begin{proof}
For any number $y$, we have
\[ F_Y(y) = \Pr(Y \le y) = \Pr(X \le (y-b)/a) = F_X((y-b)/a). \]
Hence, for any $p \in [0, 1]$, we have
\begin{align*}
Q_Y(p) \ge y &\iff F_Y(y) \le p \iff F_X((y-b)/a) \le p
\\ &\iff Q_X(p) \ge (y-b)/a \iff aQ_X(p) + b \ge y.
\end{align*}
Hence, $Q_Y(p) = aQ_X(p) + b$ for all $p \in [0, 1]$.
\end{proof}

In \cref{fig:exp-vs-norm-qq} we have the QQ-plot of the standard normal distribution
w.r.t.~the exponential distribution ($\lambda = 1$)

\begin{figure}[!htb]
\centering
\includegraphics{expVsNormQQ.pdf}
\caption{The blue line is the QQ plot of $Y \sim \Ncal(0, 1)$ w.r.t~$X \sim \Exp(1)$.
The red dashed line is the equation $y = x - 1$.}
\label{fig:exp-vs-norm-qq}
\end{figure}

\subsection{Tail Behavior from QQ plot}

The QQ plot can tell us whether two distributions are roughly the same,
but it help us in other ways too.
By analyzing how the QQ plot differs from a straight line at the left and right ends of the plot,
we can compare the tails of the two distributions.

To do that, we first draw the \emph{reference line}.
In a QQ plot of $Y$ w.r.t.~$X$, the reference line has the equation
\[ y = \E[Y] + \sqrt{\frac{\Var[Y]}{\Var[X]}}(x - \E[X]). \]

\begin{enumerate}
\item When $Y = aX + b$, the reference line becomes $y = ax + b$.
\item For the example in \cref{fig:exp-vs-norm-qq}, the reference line
    has the equation $y = x-1$ and is drawn dashed red.
\end{enumerate}

In \cref{fig:exp-vs-norm-qq}, on the right side of the figure,
the QQ plot curves downwards compared to the reference line.
This means that as $p$ increases and gets closer to 1, $Q_X(p)$ grows faster than $Q_Y(p)$.
Intuitively, when a distribution has a heavy right tail, we see more extreme (larger) values,
which suggests that $Y$ has a lighter right tail than $X$.
Indeed, $X$ has PDF $f_X(x) = e^{-x}$, and $Y$ has PDF $f_Y(x) \propto e^{-x^2/2}$.
As $x$ grows large, $f_Y(x)$ becomes much smaller than $f_X(x)$,
which tells us that $Y$ has a lighter right tail than $X$.
We can see clearly in \cref{fig:exp-vs-norm}.

\begin{figure}[!htb]
\centering
\includegraphics{expVsNorm.pdf}
\caption{The blue curve is the PDF of $Y \sim \Ncal(0, 1)$.
The red curve is the PDF of $X \sim \Exp(1)$.}
\label{fig:exp-vs-norm}
\end{figure}

Let us now look at the left side of \cref{fig:exp-vs-norm-qq}.
The QQ plot curves downwards compared to the reference line.
This means that as $p$ decreases and gets closer to 0, $Q_Y(p)$ decreases much faster than $Q_X(p)$.
Intuitively, when a distribution has a heavy left tail, we see more extreme (smaller) values,
which suggests that $Y$ has a heavier left tail than $X$.
Indeed, for $x < 0$, the PDF of $X$ is 0, and $Y$ has PDF $f_Y(x) \propto e^{-x^2/2}$,
so $Y$ has a heavier left tail than $X$.

Since $Y$ has a heavier left tail than $X$ and a lighter right tail than $Y$,
we get that $Y$ is left skewed compared to $X$.

Let us summarize our observations.

\begin{tcolorbox}[myres]
Let $X$ and $Y$ be random variables.
In the QQ plot of $Y$ w.r.t.~$X$, compared to the reference line,
\begin{enumerate}
\item If the right end of the QQ plot
    \begin{tightenum}
    \item curves downwards, then $Y$ has a lighter right tail than $X$.
    \item curves upwards, then $Y$ has a heavier right tail than $X$.
    \end{tightenum}
\item If the left end of the QQ plot
    \begin{tightenum}
    \item curves downwards, then $Y$ has a heavier left tail than $X$.
    \item curves upwards, then $Y$ has a lighter left tail than $X$.
    \end{tightenum}
\end{enumerate}
\end{tcolorbox}

\begin{tcolorbox}[progExe]
\begin{exercise}
Let $X \sim \Ncal(0, 1)$ have the standard normal distribution.
Let $Y \sim U(-1, 1)$ be uniformly distributed between $-1$ and $1$.
\begin{enumerate}[label=(\alph*)]
\item What is the quantile function of $Y$?
\item Use \texttt{matplotlib.pyplot} to draw the QQ plot of $Y$ w.r.t.~$X$. Also plot the reference line.
    Label the $x$ and $y$ axes of your plot (similar to \cref{fig:exp-vs-norm-qq}).
    \textcolor{progExeSec}{You can use \texttt{\href{https://docs.scipy.org/doc/scipy/reference/generated/scipy.stats.norm.html}{%
scipy.stats.norm}.ppf(x)} as the quantile function of $X$.}
\item From the QQ plot, what can you say about the tails of $Y$
    compared to the standard normal distribution and why?
\end{enumerate}
\end{exercise}
\end{tcolorbox}

\subsection{QQ plots of sampled data}

Often, we don't know the distribution of a random variable, but we can observe samples from it,
and use them to get an \emph{estimated QQ plot}.

Suppose we observe $n$ samples from a distribution $\Dcal$.
After sorting them in non-decreasing order, let
$(y_1, y_2, \ldots, y_n)$ be the resulting sequence.
We define the \emph{estimated quantile function} $\Qhat$ as follows:
\begin{enumerate}
\item Let $p_i \defeq \frac{i-1/2}{n}$ for all $i$ from 1 to $n$.
\item Let $\Qhat(p_i) = y_i$ for all $i$ from 1 to $n$.
\end{enumerate}

Then we draw the estimated QQ plot w.r.t.~a random variable $X$ as follows:
\begin{enumerate}
\item Let $x_i \defeq Q_X(p_i)$ for all $i$ from 1 to $n$.
\item Plot the points $(x_i, y_i)$ for all $i$ from 1 to $n$.
\end{enumerate}

We demonstrate this process in \cref{fig:exp-vs-norm-qqe}.
We sample 1000 points from the standard normal distribution,
and draw the estimated QQ plot w.r.t.~the exponential distribution with parameter 1.

\begin{figure}[!htb]
\centering
\includegraphics{expVsNormQQE.pdf}
\caption{Estimated QQ plot of 1000 samples drawn from the standard normal distribution w.r.t.~$\Exp(1)$.
Observe the similarity with \cref{fig:exp-vs-norm-qq}.}
\label{fig:exp-vs-norm-qqe}
\end{figure}

Let us try to see if the Pareto distribution follows the Central limit theorem.
To do so, we sample 1000 points from $\mathrm{Pareto}(3/2)$ and compute the average.
We repeat this process 1000 times to obtain 1000 sample averages.
If the central limit theorem holds here, then these sample averages would have a normal distribution,
and a QQ plot w.r.t.~the standard normal distribution would give us a straight line.
However, this does not happen, as seen in \cref{fig:pareto-mean-qqe}.
Hence, the central limit theorem does not hold for $\mathrm{Pareto}(3/2)$.

\begin{figure}[!htb]
\centering
\includegraphics{paretoMeanQQE.pdf}
\caption{Estimated QQ plot of 1000 sample averages of 1000 points drawn from $\mathrm{Pareto}(3/2)$
w.r.t.~$\Ncal(0, 1)$. Note how the curve quickly starts jumping up as we move to the right.
In fact, we had to remove the rightmost point $(3.2905, 259.78)$ from the plot to make it easier to see.}
\label{fig:pareto-mean-qqe}
\end{figure}

\begin{tcolorbox}[mathExe]
\begin{exercise}
Why does the central limit theorem not hold for $\mathrm{Pareto}(3/2)$?
\\ \textcolor{mathExeSec}{(Hint: recall the conditions for applying the central limit theorem.)}
\end{exercise}
\end{tcolorbox}

\bibPrep{}  % \bibPrep is defined in header.tex
\bibliography{src/bibdb}

\cite{hill1975simple}

\end{document}
